{\renewcommand{\lstlistingname}{Script}
\begin{figure*}[t]
\centering
\begin{lstlisting}[style=python,
    caption={consenso\_prom.launch.py -- Parte 2: construccion dinamica de acciones.},
    label={lst:consenso_part2}]
def launch_setup(context, *args, **kwargs):

    # Leer y convertir los argumentos ya resueltos por el contexto
    n              = int(LaunchConfiguration('n_robots').perform(context))
    R              = float(LaunchConfiguration('neighbor_radius').perform(context))
    t_final        = float(LaunchConfiguration('t_final').perform(context))
    spawn_obstacle = LaunchConfiguration('spawn_obstacle').perform(context).lower() == 'true'
    obs_x          = float(LaunchConfiguration('obstacle_x').perform(context))
    obs_y          = float(LaunchConfiguration('obstacle_y').perform(context))

    poses         = POSES_8[:n]
    simulador_pkg = get_package_share_directory('simulador')
    obstacle_sdf  = os.path.join(simulador_pkg, 'models', 'obstacle_cylinder.sdf')
    articubot_pkg = get_package_share_directory('articubot_one')
    gazebo_launch = os.path.join(get_package_share_directory('gazebo_ros'),
                                 'launch', 'gazebo.launch.py')
    rsp_launch    = os.path.join(articubot_pkg, 'launch', 'rsp.launch.py')

    actions = []

    # 1. Gazebo arranca de inmediato para tener el mayor tiempo de
    #    inicializacion antes de que lleguen los primeros spawns.
    actions.append(
        IncludeLaunchDescription(PythonLaunchDescriptionSource(gazebo_launch))
    )

    # 2. Obstaculo opcional: se spawnea a los 3 s si el argumento lo pide.
    if spawn_obstacle:
        actions.append(TimerAction(period=3.0, actions=[
            Node(package='gazebo_ros', executable='spawn_entity.py',
                 arguments=['-file', obstacle_sdf, '-entity', 'obstacle_cylinder',
                             '-x', str(obs_x), '-y', str(obs_y), '-z', '0.5'],
                 output='screen')
        ]))

    # 3. Por cada robot: RSP -> spawn -> control, con 1.5 s entre robots
    #    para no saturar Gazebo durante el arranque masivo.
    for i, (x0, y0, yaw) in enumerate(poses):
        ns    = f'robot{i}'
        delay = 3.0 + i * 1.5

        rsp_node = IncludeLaunchDescription(
            PythonLaunchDescriptionSource(rsp_launch),
            launch_arguments={'namespace': ns, 'use_sim_time': 'false'}.items()
        )
        spawn_node = Node(
            package='gazebo_ros', executable='spawn_entity.py',
            arguments=['-topic', f'/{ns}/robot_description', '-entity', ns,
                       '-x', str(round(x0,4)), '-y', str(round(y0,4)),
                       '-z', '0.1', '-Y', str(round(yaw,4))],
            output='screen'
        )
\end{lstlisting}
\end{figure*}
}
